% Part: normal-modal-logic
% Chapter: syntax-and-semantics
% Section: language-modal-logic

\documentclass[../../../include/open-logic-section]{subfiles}

\begin{document}

\olfileid{nml}{syn}{lan}

\olsection{મૂળભૂત મોડલ તર્કશાસ્ત્રની ભાષા}

\begin{defn}
મોડલ તર્કશાસ્ત્રની મૂળભૂત ભાષામાં નીચેની વસ્તુઓ હોય છે:
\begin{enumerate}
  \tagitem{prvFalse}{!!{falsity} માટેનો વિધાનાત્મક અચલ~$\lfalse$.}{}
  \tagitem{prvTrue}{!!{truth} માટેનો વિધાનાત્મક અચલ~$\ltrue$.}{}
  \item !!{propositional variable}sનો !!{denumerable}s ગણ: $\Obj
    p_0$, $\Obj p_1$, $\Obj p_2$, \dots
  \item વિધાનાત્મક સંયોજકો: \startycommalist
  \iftag{prvNot}{\ycomma $\lnot$ (નિષેધ)}{}%
  \iftag{prvAnd}{\ycomma $\land$ (સંયોજન)}{}%
  \iftag{prvOr}{\ycomma $\lor$ (વિયોજન)}{}%
  \iftag{prvIf}{\ycomma $\lif$ (!!{conditional})}{}%
  \iftag{prvIff}{\ycomma $\liff$ (!!{biconditional})}{}.
  \tagitem{prvBox}{મોડલ કારક $\Box$.}{}
  \tagitem{prvDiamond}{મોડલ કારક $\Diamond$.}{}
\end{enumerate}
\end{defn}

\begin{defn}
મૂળભૂત મોડલ ભાષાનાં \emph{!!^{formula}s} નીચે પ્રમાણે
આગમનાત્મક રીતે વ્યાખ્યાયિત થાય છે:
\begin{enumerate}
\tagitem{prvFalse}{$\lfalse$ એ પરમાણુ !!{formula} છે.}{}

\tagitem{prvTrue}{$\ltrue$ એ પરમાણુ !!{formula} છે.}{}

\item દરેક !!{propositional variable} $\Obj p_i$ એ (પરમાણુ) !!{formula} છે.

\tagitem{prvNot}{જો $!A$ એ !!a{formula} હોય, તો $\lnot !A$ પણ
  !!a{formula} છે.}{}

\tagitem{prvAnd}{જો $!A$ અને $!B$ એ !!{formula}s હોય, તો $(!A \land
  !B)$ એ !!a{formula} છે.}{}

\tagitem{prvOr}{જો $!A$ અને $!B$ એ !!{formula}s હોય, તો $(!A \lor !B)$
  એ !!a{formula} છે.}{}

\tagitem{prvIf}{જો $!A$ અને $!B$ એ !!{formula}s હોય, તો $(!A \lif !B)$
  એ !!a{formula} છે.}{}

\tagitem{prvIff}{જો $!A$ અને $!B$ એ !!{formula}s હોય, તો $(!A \liff !B)$
  એ !!a{formula} છે.}{}

\tagitem{prvBox}{જો $!A$ એ !!a{formula} હોય, તો $\Box !A$ પણ
  !!a{formula} છે.}{}

\tagitem{prvDiamond}{જો $!A$ એ !!a{formula} હોય, તો $\Diamond !A$ પણ
  !!a{formula} છે.}{}

\tagitem{limitClause}{બીજું કંઈ !!a{formula} નથી.}{}
\end{enumerate}
\end{defn}

\begin{tagblock}{defNot,defOr,defAnd,defIf,defIff,defTrue,defFalse,defBox,defDiamond}
\begin{defn}
વ્યાખ્યાયિત કારકો વડે રચાયેલાં સૂત્રોનો અર્થ નીચે મુજબ લેવો:

\begin{tagenumerate}{defTrue,defFalse,defNot,defOr,defAnd,defIf,defIff,defBox,defDiamond}

\tagitem{defTrue}{$\ltrue$ એ
  \iftag{prvFalse}{$\lnot\lfalse$}{$(!A \lor \lnot !A)$, જ્યાં
    $!A$ કોઈ નિશ્ચિત પરમાણુ !!{formula} છે}નું સંક્ષેપ છે.}{}

\tagitem{defFalse}{$\lfalse$ એ
  \iftag{prvTrue}{$\lnot\ltrue$}{$(!A \land \lnot !A)$, જ્યાં
    $!A$ કોઈ નિશ્ચિત પરમાણુ !!{formula} છે}નું સંક્ષેપ છે.}{}

\tagitem{defNot}{$\lnot !A$ એ $!A \lif \lfalse$નું સંક્ષેપ છે.}{}

\tagitem{defOr}{$!A \lor !B$ એ
  \iftag{prvAnd}{$\lnot(\lnot !A \land \lnot !B)$}{$\lnot !A \lif
    !B$}નું સંક્ષેપ છે.}{}

\tagitem{defAnd}{$!A \land !B$ એ
  \iftag{prvOr}{$\lnot(\lnot !A \lor \lnot !B)$}{$\lnot (!A \lif
    \lnot !B)$}નું સંક્ષેપ છે.}{}

\tagitem{defIf}{$!A \lif !B$ એ
  \iftag{prvOr}{$\lnot !A \lor !B$}{$\lnot (!A \land \lnot !B)$}નું સંક્ષેપ છે.}{}

\tagitem{defIff}{$!A \liff !B$ એ $(!A \lif !B) \land (!B
  \lif !A)$નું સંક્ષેપ છે.}{}

\tagitem{defBox}{$\Box !A$ એ $\lnot\Diamond\lnot !A$નું સંક્ષેપ છે. }{}

\tagitem{defDiamond}{$\Diamond !A$ એ $\lnot\Box\lnot !A$નું સંક્ષેપ છે.}{}
\end{tagenumerate}
\end{defn}
\end{tagblock}
\sourcecorrection{OLMD-001}{સ્થિર અંગ્રેજી મૂળમાં
\texttt{defIf}ની પ્રથમ શાખામાં $\lnot !A \lor !B)$
લખતાં બંધ કૌંસ માટે ખુલતું કૌંસ નથી. અહીં વધારાનું
બંધ કૌંસ દૂર કરીને $\lnot !A \lor !B$ રાખ્યું છે.}

જો !!a{formula}~$!A$માં $\Box$ કે $\Diamond$ ન હોય,
તો તેને \emph{મોડલ-મુક્ત} કહીએ છીએ.

\end{document}
